% ICASSP 2027 -- 4 pages technical + page 5 references only.
% EVERY number here must be reproducible by code/verify_paper.py from results/*.json.
% No number from the lost pre-2026-09-11 runs.
\documentclass{article}
\usepackage{spconf,amsmath,amssymb,graphicx,booktabs,url}
\usepackage[hidelinks]{hyperref}

\title{Cosine-Similarity Trust Constrains One Direction:\\ A Closed-Form Subspace Backdoor Against Norm-Equalized Robust Aggregation}
%
% ICASSP single-blind. Zhou: three signing affiliations (confirmed 2026-09-05);
% Dai, Sha, Yang: School of Computer Science, NJUPT. Corresponding: Dai, Sha.
\name{Hao Zhou$^{1,2,3}$, Hua Dai$^{1,\dagger}$, Letian Sha$^{1,\dagger}$, Geng Yang$^{1}$%
\thanks{$^{\dagger}$Corresponding authors. This work was supported in part by the
National Natural Science Foundation of China under Grant~62602293, in part by the
Natural Science Foundation of Jiangsu Province under Grant~BK20250683, and in part
by the Jiangsu Province Youth Science and Technology Talent Support Project under
Grant~JSTJ-2025-641.}}
\address{$^{1}$School of Computer Science, Nanjing University of Posts and Telecommunications, Nanjing, China\\
  $^{2}$State Key Laboratory for Novel Software Technology, Nanjing University, Nanjing, China\\
  $^{3}$Tongda College of Nanjing University of Posts and Telecommunications, Yangzhou, China}

\begin{document}
\maketitle

\begin{abstract}
Secure aggregation hides individual client updates from the server, while poisoning
defenses need to inspect them. Similarity-based trust aggregation (FLTrust, RFLPA)
reconciles the two by scoring each client with a single scalar --- the cosine
between its update and a private server-side reference --- and equalizing update
norms. We show that this pair of guards fixes the \emph{magnitude} of the
component orthogonal to the reference, $\|g_0\|\sqrt{1-\tau^{2}}$ at cosine
$\tau$, and leaves its \emph{direction} free. A backdoor can therefore be written
in closed form --- a cover along the adversary's estimate of the reference plus a
payload in its orthogonal complement --- with no constraint-satisfaction search
and no observation of honest updates. At $30\%$ malicious clients on Fashion-MNIST, a \emph{naive} backdoor already
leaks through this defense --- but only intermittently, succeeding on $3$ of $5$
seeds with no trust advantage ($1.04\times$ the honest score). Aiming the cover at the estimated anchor makes the attack reliable: our construction succeeds in all five tested seeds ($0.982\pm0.014$) by earning a $1.44\times$ trust advantage, against $3/5$ for the naive attempt. A random-direction cover, whose trust falls far below the honest score, succeeds in $0/5$ at $\tau=0.97$. The construction does not break an otherwise-robust defense; it turns an unreliable leak into a consistent one, and we identify the mechanism --- the anchor-aligned cover --- that does it. We will
release the harness and all raw results.%
% CAMERA-READY: make the repo public and uncomment the next line (drop the leading %),
% then remove the "will" above so it reads "We release ...":
% \footnote{\url{https://github.com/haozhou12/cpsubspace-fl-backdoor}}
\end{abstract}

\begin{keywords}
Federated learning, robust aggregation, secure aggregation, backdoor, cosine similarity
\end{keywords}

% ====================================================================
\section{Introduction}
\label{sec:intro}
Federated learning (FL) is expected to give two guarantees at once. Secure
aggregation (SecAgg)~\cite{secagg} reveals only the sum of client updates to the
server; robust aggregation defends the global model against clients that submit
poisoned updates. The two pull in opposite directions: a poisoning defense must
score individual updates, and SecAgg exists to make them invisible.

The dominant reconciliation is \emph{similarity-based trust}. FLTrust~\cite{fltrust}
defines it: the server keeps a small clean root set, computes a reference update
$g_0$ from it, scores each client by the cosine between its update and $g_0$,
normalizes updates to $\|g_0\|$, and takes the trust-weighted mean.
RFLPA~\cite{rflpa} implements the same rule under secure aggregation, evaluating
the score through a privacy-preserving dot-product protocol so the server never
sees an individual update, and keeps $g_0$ private --- a client that cannot see
the reference should not be able to aim at it.

This paper shows that the statistic is weaker than that intuition suggests, in
two stages. First, structurally: a cosine score constrains an update along one
direction; on the $(d\!-\!1)$-dimensional orthogonal complement the score is
invariant, and norm equalization fixes magnitude, not direction. Second,
empirically: at a high malicious fraction the defense is already porous --- a
plain poisoned update slips through whenever its direction happens to correlate
with the reference --- but unreliably. A closed-form cover aimed at the estimated reference lifts the adversary's trust above the honest score and plants the backdoor consistently across seeds. Related high-dimensional freedom has
been used against Byzantine-robust aggregators~\cite{hidra} and by perturbations
that stay inside the benign variance~\cite{alie}; neither targets a private-anchor
trust score. Our contributions are:
\begin{itemize}\setlength\itemsep{0pt}
  \item[(i)] a characterization of the feasible set left by cosine trust with
    norm equalization --- a sphere of dimension $d-2$ on which the defense cannot
    distinguish points --- and of the adversary's payload budget on it
    (Sec.~\ref{sec:attack});
  \item[(ii)] a closed-form, norm-matched backdoor construction that needs no
    constraint-satisfaction search and no honest update, with an exact gray-box
    trust band (Sec.~\ref{sec:attack});
  \item[(iii)] an evaluation against the \emph{published} unconditional normalization rule
    showing that the anchor-aligned cover lifts an intermittent naive backdoor ($3/5$ seeds) to all five tested seeds by earning a $1.44\times$ trust advantage, that a random cover earns none and succeeds $0/5$ at $\tau{=}0.97$, and a $\tau$ sweep exhibiting two failure modes consistent with the analysis
    (Sec.~\ref{sec:exp}).
\end{itemize}
We do not compare against optimization-based adaptive attacks~\cite{autoadapt}:
we did not locate a public implementation, and a faithful reimplementation is
outside the scope of this letter. Our claim is not that the attack is stronger than prior poisoning, but that the geometry makes it \emph{closed-form and reliable} where a naive attempt is a coin flip.

% ====================================================================
\section{System and Threat Model}
\label{sec:model}
\noindent\textbf{Aggregation.} $n$ clients; in round $t$ client $i$ submits
$g_i\in\mathbb{R}^d$ under SecAgg~\cite{secagg}. The server computes a reference
$g_0$ from its private root set and, inside the protocol~\cite{rflpa}, assigns
\begin{equation}
  \mathrm{TS}_i=\max\!\bigl(0,\cos(g_i,g_0)\bigr),\qquad
  \bar g_i=\frac{\|g_0\|}{\|g_i\|}\,g_i ,
  \label{eq:ts}
\end{equation}
i.e.\ every update is normalized to the reference norm (RFLPA Eq.~3), and takes
$G=\sum_i\mathrm{TS}_i\bar g_i/\sum_i\mathrm{TS}_i$. We also evaluate the
one-sided variant $\bar g_i=g_i\min(1,\|g_0\|/\|g_i\|)$, which leaves updates
shorter than $\|g_0\|$ unchanged. As $\cos$ is scale-invariant the two rules give
identical trust; they differ only in the mass a short honest update contributes.
$g_0$ is never released; RFLPA does distribute its norm $\|g_0\|$.

\noindent\textbf{Adversary (gray-box).} A fixed fraction $\rho$ of clients
($30\%$ unless stated) is malicious and may coordinate. They know \eqref{eq:ts},
the current global model, and $\|g_0\|$. They do \emph{not} know the reference
direction, the root set, or any honest update --- SecAgg denies the latter two.
They substitute an estimate $\hat e_0$ of $g_0/\|g_0\|$ computed from their own
data. The goal is a backdoor: inputs bearing a trigger are classified into a
target label while clean accuracy is preserved. The server is honest; privacy
attacks on SecAgg~\cite{graddisagg,eluding} are out of scope.

% ====================================================================
\section{Closed-Form Subspace Construction}
\label{sec:attack}

\subsection{What the two guards leave free}
For \emph{any} update with $\|g\|=R$ and $\cos(g,a)=\tau$ against a reference $a$,
the component along $a$ is $\tau R$ and the residual orthogonal to $a$ has norm
\begin{equation}
  \bigl\|g-\langle g,\hat a\rangle\hat a\bigr\| = R\sqrt{1-\tau^{2}},
  \label{eq:budget}
\end{equation}
independently of how the residual is oriented. Trust and normalization therefore
fix the payload's \emph{magnitude} and leave its \emph{direction} free: at given
$(R,\tau)$ the feasible set is a sphere of dimension $d-2$, and with respect to
the server's own anchor every point on it is equally acceptable. The gray-box
case, where the adversary can only aim at an estimate of that anchor, is treated
in Sec.~\ref{sec:graybox}.

\subsection{Construction}
Fig.~\ref{fig:geom} summarizes the geometry. Let $b$ be a local update on the
poisoned objective. Remove its anchor component and normalize,
\begin{equation}
  v=\frac{b-\langle b,\hat e_0\rangle\hat e_0}{\|b-\langle b,\hat e_0\rangle\hat e_0\|},
  \qquad \langle v,\hat e_0\rangle=0,\ \|v\|=1,
  \label{eq:payload}
\end{equation}
and for a scalar $\tau\in(0,1)$ let every colluding client submit
\begin{equation}
  g_A=\|g_0\|\bigl(\tau\,\hat e_0+\sqrt{1-\tau^{2}}\;v\bigr),
  \label{eq:attack}
\end{equation}
the \emph{cover} and the \emph{payload}. Since $\hat e_0\perp v$ are unit vectors,
$\|g_A\|=\|g_0\|$ exactly, so under either normalization rule $\bar g_A=g_A$ and
the payload arrives with its full magnitude $\|g_0\|\sqrt{1-\tau^{2}}$. Under the
unconditional rule the submitted scale is immaterial, since every update is
rescaled to $\|g_0\|$; under the one-sided clip any scale $\ge\|g_0\|$ gives the
same result. Under a perfect anchor $\cos(g_A,g_0)=\tau$, so $\tau$ is the trust
score itself; and among all $g$ with $\|g\|=\|g_0\|$ and trust $\tau$,
\eqref{eq:attack} is the one that spends the entire budget \eqref{eq:budget} on
the backdoor direction. Each round costs the local training that produces $b$ (in our implementation, ten SGD steps to estimate $\hat e_0$ and thirty on the poisoned objective) and one $O(d)$ projection; there is no loop to satisfy the trust and norm constraints and no server feedback, which is the structural difference from adaptive poisoning.

\subsection{Gray-box trust band}
\label{sec:graybox}
In practice $\hat e_0\neq\hat g_0\triangleq g_0/\|g_0\|$. With
$\delta=\angle(\hat e_0,\hat g_0)$,
\begin{equation}
  \cos(g_A,g_0)=\underbrace{\tau\cos\delta}_{\text{controlled}}
  +\underbrace{\sqrt{1-\tau^{2}}\,\langle v,\hat g_0\rangle}_{\text{residual}},
  \label{eq:true}
\end{equation}
where the residual is not chosen by the adversary: $v\perp\hat e_0$ but not
$\perp g_0$. Since $|\langle v,\hat g_0\rangle|\le\sin\delta$, the true cosine
lies between $\tau\cos\delta\mp\sqrt{1-\tau^{2}}\sin\delta$, i.e.
\begin{equation}
  \cos(\alpha+\delta)\le\cos(g_A,g_0)\le\cos(\alpha-\delta),\quad\alpha=\arccos\tau,
  \label{eq:band}
\end{equation}
tight when $v\in\mathrm{span}(\hat e_0,\hat g_0)$. Two consequences follow.
First, the clamp in \eqref{eq:ts} zeroes the adversary once the sum in
\eqref{eq:true} is nonpositive; the controlled term $\tau\cos\delta$ stays
positive while the anchor estimate beats chance, and shrinks with a poor
estimate. Second, the lever the adversary holds is the payload budget
$\sqrt{1-\tau^{2}}$, which \emph{decreases} with $\tau$: raising $\tau$ trades
payload for trust. Neither limit is the whole story --- Sec.~\ref{sec:exp} sweeps
$\tau$ and finds a seed failing by clamping at low $\tau$ and every seed failing
by starvation as $\sqrt{1-\tau^{2}}\!\to\!0$.

\begin{figure}[t]
  \centering
  \includegraphics[width=\columnwidth]{figs/fig1_geometry.pdf}
  \caption{Geometry of \eqref{eq:attack}. The cover fixes the cosine to the
  \emph{estimated} anchor at $\tau$; the payload rides in the orthogonal
  complement with magnitude $\|g_0\|\sqrt{1-\tau^{2}}$; $\|g_A\|=\|g_0\|$ so
  normalization is the identity on $g_A$.}
  \label{fig:geom}
\end{figure}

% ====================================================================
\section{Experiments}
\label{sec:exp}

\noindent\textbf{Setup.} Fashion-MNIST, a two-convolution CNN with
GroupNorm~\cite{groupnorm} (BatchNorm statistics are buffers and are not
aggregated under update-based FL), $20$ clients with an IID split disjoint from a
$100$-sample root set, $30\%$ malicious, $30$ local SGD steps per round,
$120$ rounds. The adversary estimates $\hat e_0$ with $10$ local steps on its own
data. Trigger: a $3\times3$ patch, target class~$0$; ASR is measured on triggered
\emph{non-target} test inputs. Unless stated, $\tau=0.97$ and the published
unconditional rule \eqref{eq:ts} is used. We evaluate the unquantized aggregation rule \eqref{eq:ts} in a centralized harness and do not implement RFLPA's cryptographic protocol; FLTrust defines the rule, RFLPA realizes it securely. Numbers are means $\pm$ population s.d.\ over seeds, with per-seed values released; we never report a maximum over
seeds. GPU training is not bitwise deterministic, so re-running a seed can move
final ASR by a few points; success counts for cells near the $0.8$ threshold are
therefore indicative.

\noindent\textbf{Main result: the defense is porous, the cover makes it
reliable (Table~\ref{tab:main}, Fig.~\ref{fig:mech}).} A \emph{naive} backdoor
--- the poisoned local update submitted directly, no cover, no projection ---
already leaks through at $30\%$ malicious clients, but only when its own direction happens to earn positive trust: it succeeds in exactly the $3/5$ seeds whose update clears the cosine clamp (mean $\overline{\mathrm{TS}}_{\mathrm{mal}}=0.076$ for the successes versus $0.004$ for the two failures), a bimodal outcome ($0.675\pm0.399$; per-seed ASR $0.05$ to $1.0$) with no trust advantage ($1.04\times$ the honest score). The anchor-aligned cover removes this dependence on chance: it lifts the malicious trust to $1.44\times$ and plants the backdoor in all five tested seeds ($0.982\pm0.014$). The advantage comes from the alignment, not from submitting a norm-matched update at all: a random-direction cover earns trust far below the honest score ($0.09\times$) and succeeds in $0/5$. The result is rule-invariant: under the one-sided clip the
construction reaches $0.996\pm0.003$ (paired difference $-0.014$, $p=0.073$). It
costs $2.6$ clean-accuracy points at this $\tau$ ($0.861$ vs $0.887$ attack-free,
Welch $p=0.013$).

\begin{table}[t]
  \centering\small
  \caption{Fashion-MNIST, 120 rounds, 30\% malicious. Final ASR / clean accuracy
  (mean $\pm$ s.d.); succ.\ = seeds with ASR $\ge0.8$; trust = malicious over honest mean $\mathrm{TS}$ (each averaged over seeds and the last 30 rounds; each row uses its own honest scores). Norm.\ = published rule \eqref{eq:ts}; clip = one-sided
  variant. All rows $n{=}5$ except none ($n{=}3$).}
  \label{tab:main}
  \setlength{\tabcolsep}{3pt}
  \begin{tabular}{@{}l@{\hspace{4pt}}cccc@{}}
    \toprule
    Rule / attack & ASR & succ. & clean & trust \\
    \midrule
    norm. / none                & $0.007\pm0.003$ & 0/3 & $0.887$ & --- \\
    norm. / naive backdoor      & $0.675\pm0.399$ & 3/5 & $0.875$ & $1.04\times$ \\
    norm. / random cover        & $0.029\pm0.037$ & 0/5 & $0.881$ & $0.09\times$ \\
    norm. / ours $\tau{=}0.97$  & $\mathbf{0.982\pm0.014}$ & 5/5 & $0.861$ & $\mathbf{1.44\times}$ \\
    clip / ours $\tau{=}0.97$   & $0.996\pm0.003$ & 5/5 & $0.857$ & --- \\
    norm. / ours $\tau{=}0.7$   & $0.977\pm0.046$ & 5/5 & $0.878$ & $1.43\times$ \\
    \bottomrule
  \end{tabular}
\end{table}

\begin{figure}[t]
  \centering
  \includegraphics[width=\columnwidth]{figs/fig_mech.pdf}
  \caption{Per-seed final ASR at $30\%$ malicious (bars are means). A naive
  backdoor is bimodal --- a coin flip; a random-direction cover fails; the
  anchor-aligned cover succeeds in every seed.}
  \label{fig:mech}
\end{figure}

\noindent\textbf{Two failure modes (Fig.~\ref{fig:tau}).} Sweeping $\tau$ at
$30\%$ malicious clients ($n{=}5$ per point), the attack succeeds in every seed
at $\tau\in\{0.7,0.85,0.97\}$ and collapses to $0.492\pm0.144$ ($0/5$) at
$\tau=0.995$, where the payload budget $\sqrt{1-\tau^{2}}$ is $0.10$; the
adversary's mean trust is still positive there
($\mathrm{TS}_{\mathrm{mal}}=0.21$), so this end is payload starvation, not the
clamp. Below $0.7$, \emph{one} seed of five fails at both $\tau=0.5$ and
$\tau=0.3$ (ASR $0.158$, $0.051$); it has the poorest anchor estimate of the
five ($\cos(\hat e_0,g_0)$ averages $-0.07$ over the last thirty rounds at
$\tau=0.5$, versus $+0.01$ to $+0.16$ for the others), so its controlled term
$\tau\cos\delta$ cannot keep the sum positive and \eqref{eq:ts} clamps it to
$0$ in over $90\%$ of rounds; at $\tau=0.7$ the same seed is clamped in $73\%$ of
rounds and still reaches ASR $0.886$. Both ends are consistent with the limits of \eqref{eq:true}; we do not claim a universal sweet spot. The adversary's mean per-round
aggregate share (its fraction of $\sum_i\mathrm{TS}_i\|\bar g_i\|$, last thirty rounds) is $0.34$--$0.38$ over $\tau\in[0.7,0.995]$ but falls to $0.25$ at
$\tau=0.5$ and $0.17$ at $\tau=0.3$, tracking the clamping just described.

\begin{figure}[t]
  \centering
  \includegraphics[width=\columnwidth]{figs/fig3_tau.pdf}
  \caption{$\tau$ sweep (Fashion-MNIST, 120 rounds, 30\% malicious, published
  rule; $n{=}5$, dots are individual seeds). Low $\tau$ fails by clamping in one
  seed; high $\tau$ fails every seed by payload starvation.}
  \label{fig:tau}
\end{figure}

\noindent\textbf{Where the clean cost comes from.} The cost is not monotone in $\tau$ (it rises to $2.8$ points at $\tau=0.85$ then falls). Two component swaps place it with the cover, not the backdoor. A random orthogonal payload drives ASR to the no-attack level yet still costs $2.1$ points at $\tau=0.97$ and $0.9$ at $\tau=0.7$; a cover-only update ($\tau=1$, no payload) costs $1.4$, against $2.6$ for the full attack. Clean loss thus occurs without a successful backdoor; the payload direction is not what the operator pays for. We do not further separate the cover and payload contributions.

\noindent\textbf{Malicious fraction.} The naive backdoor stays a lottery as the
fraction drops --- $1/5$ at $10\%$ and $3/5$ at $20\%$, always bimodal --- and our construction likewise thins out as the fraction drops: at $\tau=0.7$ it reaches $4/5$ at $20\%$ ($0.747\pm0.365$) and $0/5$ at $10\%$ ($0.410\pm0.293$), and at $\tau=0.97$ it is $1/3$ and $0/3$. At $30\%$ the $1.44\times$ trust edge raises the malicious weight share to $0.39$; reliable planting at $10\%$ is beyond both attacks in this setting.

\noindent\textbf{Onset.} ASR lags convergence: mean clean accuracy is within
$0.05$ of its final value from round~$25$ (where mean ASR is still $0.20$), and
individual seeds first reach ASR $0.5$ at rounds $25$ to $44$.
A separate $40$-round budget yields $0.512\pm0.267$. We make no claim its onset
can be predicted.

\noindent\textbf{Cross-dataset check.} On CIFAR-10 with the \emph{same} compact
CNN (a weak model here: attack-free clean accuracy $0.550\pm0.024$, $n{=}5$), the
construction reaches ASR $0.946\pm0.023$ ($5/5$) at $\tau=0.7$ and
$0.831\pm0.035$ ($2/3$, $n{=}3$) at $\tau=0.97$, against $0.040$ attack-free.
The CIFAR-10 clean costs are $6.0$ points ($p=0.004$, $n{=}5$ vs $5$) at $\tau=0.7$ and $9.8$ ($p=0.04$, $n{=}3$ vs $5$) at $\tau=0.97$, both larger than on Fashion-MNIST, where the cost is $2.6$ points ($p=0.013$) at $\tau=0.97$ and under one point (not significant) at $\tau=0.7$. This shows the construction is not specific to one dataset; it is not a claim about large models.

\noindent\textbf{Untargeted controls.} Two classical manipulations of a
\emph{clean} local update --- scaling by $10$ and sign flipping --- give ASR
$0.008\pm0.004$ and $0.029\pm0.009$ ($n{=}3$). Neither attempts the trigger, so
low ASR is not evidence of rejection. Scaling is neutralized exactly by
\eqref{eq:ts} (clean accuracy $0.889$, versus $0.887$ attack-free). Sign flipping
is \emph{not} rejected: mean malicious trust is positive in $71$--$77\%$ of
rounds and it costs $7.2$ points of clean accuracy ($p=0.06$). It is simply not a
backdoor.

% ====================================================================
\section{Conclusion}
\label{sec:concl}
Cosine-similarity trust with norm equalization constrains one direction of a
client update and fixes the magnitude, but not the direction, of everything
orthogonal to it. A closed-form, norm-matched construction exploits exactly that:
it plants a backdoor under the published RFLPA rule with no constraint-satisfaction
search. Its clean-accuracy cost --- $2.6$ points at $\tau=0.97$ on Fashion-MNIST, smaller at $\tau=0.7$, larger on CIFAR-10 --- is incurred by the cover, not the backdoor payload: a random or absent payload costs nearly as much while planting nothing. The analysis suggests two failure mechanisms, clamping of a poor anchor estimate and payload starvation, and both appear in the sweep. A defense that scores only a scalar projection cannot see the payload;
scoring the orthogonal component, or the coherence of residuals across clients,
is the natural next step.

\vspace{2pt}\noindent\textbf{Compliance with ethical standards.} This is a
theoretical and simulation study of an attack on a published defense; it used
only public datasets (Fashion-MNIST, CIFAR-10) and involved no human subjects,
personal data, or deployed system. We study this attack to expose a structural
weakness in cosine-similarity trust so that defenses can address it, not to
enable misuse: the construction needs control of a large client fraction, degrades
clean accuracy, and --- as our results show --- is mitigated by scoring the
component orthogonal to the reference. We disclosed our findings to the authors of
the studied defense prior to publication and release our code to support
verification and the development of countermeasures.

\vfill\pagebreak
\bibliographystyle{IEEEbib}
\bibliography{refs}
\end{document}
